\section{A quadratic CM period witness with the required metric}
\label{sec:sr-quadratic-period-witness}

The flat determinant comparison uses the quadratic subfield
\[
                  K=\mathbb Q(\sqrt{-D}),
                  \qquad D=\operatorname{Host.d}.
\]
This is the coordinate selected by
\(\operatorname{Host.rt}^2=-D\).
It is not the other quadratic coordinate
\(\mathbb Q(\sqrt{-\operatorname{Host.q}})\).
Here \(D\ge7\) is squarefree and \(D\equiv3\pmod4\);
thus \(\operatorname{disc}K=-D\).
Write \(\chi=\chi_{-D}\) for the primitive odd quadratic
Dirichlet character.

We construct the elliptic period witness and compute its
normalization.  We use the first main theorem of elliptic
complex multiplication, in the form that CM \(j\)-invariants
are integral and that the Hilbert class field acts
transitively through the ideal classes.
A reference with these precise assertions is Milne,
\emph{Modular Functions and Modular Forms},
Corollary~12.16 and Theorem~12.18,
\url{https://www.jmilne.org/math/CourseNotes/MF110.pdf}.
We also use potential good reduction for an elliptic curve
with integral \(j\), as in Silverman,
\emph{The Arithmetic of Elliptic Curves}, VII.5.5,
\url{https://www.math.ens.psl.eu/~obenoist/refs/Silverman.pdf}.
The analytic mean formula is derived below.

\subsection{The integral curve and its tangent generator}

Begin with the complex elliptic curve
\(\mathbb C/\mathcal O_K\), with its canonical principal
polarization and its \(\mathcal O_K\)-action.
Complex multiplication supplies a number field of definition.
Enlarge it to contain \(K\) and to define the endomorphisms.
Its integral \(j\)-invariant implies potential good reduction
at every finite prime.  A single finite extension gives good
reduction at all primes: there are only finitely many bad
primes, one may realize the required local extensions by
approximation, and good reduction persists under further
base change.
We obtain an elliptic scheme \(\mathcal C/\mathcal O_M\)
with its canonical principal polarization.
Its generic-fiber endomorphisms extend to the scheme by the
N\'eron property of a good elliptic model.

We may further arrange that the Hodge line
\(e^*\Omega^1_{\mathcal C/\mathcal O_M}\) is free.
Indeed represent it by an invertible fractional ideal \(I\).
Choose \(h\ge1\) with \(I^h=(\alpha)\), and adjoin an
\(h\)-th root \(b\) of \(\alpha\).
Then \((I\mathcal O_{M'})^h=(b)^h\).
The fractional-ideal group of the Dedekind ring
\(\mathcal O_{M'}\) has no torsion, so
\(I\mathcal O_{M'}=(b)\).
Taking this base change and, if desired, its normal closure
makes the Hodge line free without losing good reduction.

Choose a basis \(\omega\) of the Hodge line, and let \(v\)
be its dual integral tangent generator.
Use the \(\mathcal O_K\)-action for which
\[
               w=(1+\sqrt{-D})/2
\]
acts at the distinguished complex embedding by that scalar.
The same scalar identity holds on the one-dimensional
generic Lie algebra and on its integral lattice, since
both maps are integral and agree generically.
In particular
\[
                     w^2-w+(1+D)/4=0
\]
is the derivative relation required in the period witness.
The Rosati involution is conjugation on \(\mathcal O_K\).
We also choose the representing line bundle to satisfy
the source's primitive-pairing identity, not only equality
of polarizations.  On rigidified line bundles let
\(T(L)\) be the cubical pairing pulled back by \(1,w\).
The theorem of the cube makes \(T\) additive and makes
it kill \(\operatorname{Pic}^0\).
On the N\'eron--Severi group its action is multiplication
by \(w+\overline w=1\).
Therefore \(T^2=T\): the quotient \(T(L)/L\) is in
\(\operatorname{Pic}^0\), where \(T\) is zero.
Starting with the symmetric degree-one bundle of the
zero section, choose its image under \(T\).
It is again symmetric of degree one and represents the
same principal polarization, and now
\(T(L)\simeq L\) as rigidified line bundles.
The identity extends over the good model; the relative
Picard scheme is separated and the rigidification
removes a line pulled back from the base.
This is precisely \srcid{Ell.prim}.
The metric on the tangent space depends on the
polarization, so this choice does not change its
normalization.
At the distinguished embedding the period lattice is
homothetic to \(\mathcal O_K=\mathbb Z+\mathbb Zw\),
because the chosen complex curve was
\(\mathbb C/\mathcal O_K\).
No assertion that every Galois conjugate is in the principal
ideal class is needed.

For an embedding \(\sigma:M\hookrightarrow\mathbb C\), write
the period lattice for \(\omega\) as
\[
                   \Lambda_\sigma=
                     \Omega_\sigma(\mathbb Z+\tau_\sigma\mathbb Z),
       \qquad y_\sigma=\operatorname{Im}\tau_\sigma>0 .
\]
The principal polarization has Hermitian form
\(|z|^2/\operatorname{area}(\Lambda_\sigma)\).
The source defines the tangent norm by multiplying this
form by \(\pi\).
Since \(\omega(v)=1\), its squared tangent norm is therefore
\begin{equation}\label{eq:sr-cm-tangent-norm}
                  n_\sigma=
                   \frac{\pi}{|\Omega_\sigma|^2y_\sigma}.
\end{equation}
This fixes the normalization directly, without choosing
among conventions for Faltings metrics.

The canonical discriminant section is a nowhere-zero
section of the twelfth Hodge power on a good elliptic scheme.
Relative to \(\omega^{12}\) its coefficient
\(\Delta_\omega\) is an algebraic unit.
The analytic discriminant identity is
\[
             \sigma(\Delta_\omega)
               =(2\pi)^{12}\Omega_\sigma^{-12}
                                      \eta(\tau_\sigma)^{24}.
\]
It follows either from the usual Weierstrass
\(g_2^3-27g_3^2\) identity, with invariant differential
\(dz\), or by comparing its weight-twelve Fourier
expansion at the cusp.
The product formula for the unit gives
\[
 \frac1{[M:\mathbb Q]}\sum_\sigma\log|\Omega_\sigma|
       =\log(2\pi)+
          \frac2{[M:\mathbb Q]}\sum_\sigma
                                      \log|\eta(\tau_\sigma)|.
\]
Combining this with
\eqref{eq:sr-cm-tangent-norm} yields
\begin{equation}\label{eq:sr-cm-eta-mean}
 -\frac1{[M:\mathbb Q]}\sum_\sigma\log n_\sigma
       =\frac1{[M:\mathbb Q]}\sum_\sigma
                  \log\!\left(4\pi y_\sigma
                                      |\eta(\tau_\sigma)|^4\right).
\end{equation}
CM reciprocity distributes the ideal classes uniformly
among these embeddings.
Passing to a finite extension repeats each embedding the
same number of times.  Complex conjugation preserves
\(y|\eta(\tau)|^4\).
Consequently the right side is the uniform average over
the \(h_K\) ideal classes, with \(\tau_\mathfrak a\)
any corresponding upper-half-plane class point.

\subsection{The mean from Poisson summation}

Put
\[
 \mathcal E(\tau,s)
    =\sum_{(m,n)\ne(0,0)}
                  \frac{y^s}{|m+n\tau|^{2s}}
                  \qquad(\operatorname{Re}s>1).
\]
Mellin inversion of the Gaussian and one-dimensional
Poisson summation give
\[
 \begin{split}
 \mathcal E(\tau,s)
   ={}&2\zeta(2s)y^s
       +2\sqrt\pi\,
          \frac{\Gamma(s-\tfrac12)}{\Gamma(s)}
                         \zeta(2s-1)y^{1-s}\\
     &+\frac{8\pi^s\sqrt y}{\Gamma(s)}
        \sum_{m,n\ge1}(m/n)^{s-1/2}
           K_{s-1/2}(2\pi mn y)\cos(2\pi mn\,\operatorname{Re}\tau).
 \end{split}
\]
To obtain it, separate \(n=0\), write each remaining
power as its Gaussian Mellin integral, and Poisson-sum
the \(m\)-variable.
The zero Fourier mode is the gamma quotient in the first
line, and the two opposite nonzero modes combine into
the cosine in the second.
For fixed \(y>0\), the second line converges normally
near \(s=1\), by exponential decay in \(mn y\).

At \(s=1\) use
\(K_{1/2}(t)=\sqrt{\pi/(2t)}e^{-t}\),
\(\Gamma'(\tfrac12)/\Gamma(\tfrac12)
-\Gamma'(1)/\Gamma(1)=-2\log2\),
and
\[
 \log|\eta(\tau)|
     =-\pi y/12-
         \sum_{m,n\ge1}\frac{e^{-2\pi mn y}}n
                       \cos(2\pi mn\,\operatorname{Re}\tau).
\]
The term \(2\zeta(2)y=\pi^2y/3\) cancels the
linear \(y\)-term from the eta product.
The resulting Laurent expansion is
\begin{equation}\label{eq:sr-cm-limit-formula}
 \mathcal E(\tau,s)
       =\frac{\pi}{s-1}
           +2\pi\{\gamma-\log2-
                         \log(\sqrt y|\eta(\tau)|^2)\}
           +O(s-1).
\end{equation}
All identities here use ordinary Lebesgue measure and
the Gaussian whose Fourier transform is itself.

Let \(u_K=|\mathcal O_K^\times|\).
Writing a fractional ideal as
\(\Omega(\mathbb Z+\tau\mathbb Z)\) and using
\(\operatorname{area}(\mathfrak a)
   =N(\mathfrak a)\sqrt D/2\)
shows, by grouping nonzero elements modulo units, that
\[
 \zeta_K(s)=\frac1{u_K}
           \left(\frac2{\sqrt D}\right)^s
                         \sum_{[\mathfrak a]}
                                     \mathcal E(\tau_\mathfrak a,s).
\]
This also fixes the factor \(2/\sqrt D\).
On the other hand \(\zeta_K(s)=\zeta(s)L(s,\chi)\).
Comparing residues in
\eqref{eq:sr-cm-limit-formula} gives
\[
                     L(1,\chi)=
                        \frac{2\pi h_K}{u_K\sqrt D}>0 .
\]
Comparing constant terms divided by the residue gives
\[
 \gamma+\frac{L'(1,\chi)}{L(1,\chi)}
    =2\gamma-\log2-\tfrac12\log D
       -\frac1{h_K}\sum_{[\mathfrak a]}
                   \log\!\left(y_\mathfrak a
                                  |\eta(\tau_\mathfrak a)|^4\right).
\]
Adding \(\log(4\pi)\) and using
\eqref{eq:sr-cm-eta-mean} proves the exact period value
\begin{equation}\label{eq:sr-cm-period-exact}
 -\frac1{[M:\mathbb Q]}\sum_\sigma\log n_\sigma
       =x_D
       :=\gamma+\log(2\pi/\sqrt D)
                            -\frac{L'(1,\chi)}{L(1,\chi)} .
\end{equation}
The value is real, since \(\chi\) is real and
\(L(1,\chi)>0\).

Thus the constructed integral elliptic curve with
tangent generator satisfies precisely
\(\srcid{Ell.HasPeriod}(x_D)\).
Conjugating its \(\mathcal O_K\)-action gives the witness
for the other root of \(-D\), with the same underlying
polarized curve and the same mean norm.
The construction can be based-changed to contain any
fixed number field needed by the host; the normalized
mean in \eqref{eq:sr-cm-period-exact} is preserved.

\subsection{The host discriminant correction}

For the larger CM host \(E\), the source defines
\[
 \operatorname{logOm}(E,\chi)
    =\gamma+\log(2\pi/\operatorname{rd}E)
                            -\frac{L'(1,\chi)}{L(1,\chi)} .
\]
Consequently the precise relation between its parameter
and the quadratic elliptic period is
\begin{equation}\label{eq:sr-cm-period-host}
 \operatorname{Re}\operatorname{logOm}(E,\chi)
       =x_D-\log\operatorname{rd}E+\log\sqrt D .
\end{equation}
Replacing \(\sqrt D\) by \(\operatorname{rd}E\) inside
the elliptic formula would lose this correction.

There is an equally concrete finite-place correction
when passing from the raw flat-frame norm to the
normalized determinant ratio.
For an embedding \(E\hookrightarrow M\), let
\[
 D_v=\log\|\mathfrak D_{E/\mathbb Q}\|_v
\]
use the source's ideal-log convention, in which the
normalized sum over finite places and embeddings is
\(-\log\operatorname{rd}E\).
The product formula gives the normalized finite sum of
\(\log|D|_v\) as \(-\log D\).
Thus the correction
\[
                       q(D_v-\log|D|_v)
\]
has normalized finite sum
\(q(\log D-\log\operatorname{rd}E)\).
Adding it to the finite log norm subtracts this quantity
from its Arakelov degree.
In particular the raw degree
\[
 \frac{p-q}{2}\operatorname{Re}\operatorname{logOm}(E,\chi)
          +q(\log D-\log\operatorname{rd}E)
\]
becomes exactly
\[
                 \frac{p-q}{2}
                       \operatorname{Re}\operatorname{logOm}(E,\chi).
\]
This is the correction in
\srcid{KroneckerCharacterPeriodFormula.degree\_frames}.
The calculation checks its sign and normalizes the
period witnesses just constructed.
It does not replace the separate assertion that a
given global flat frame has these integral local norms,
nor does it construct the root tower required for
\srcid{CovolumeBoundOfX}.
